\scanpage{118}

% manuscript lines / position: 1–2
i.e., en termes de $\unsure{\mathscr X|\mathcal X}=\operatorname{Spec} A$, le composé

% manuscript lines / position: 3
\[
 \mathscr X\unsure{\mathord{-}|\longrightarrow} X_1\times_X X_2\times_X X_3\longrightarrow X_1\wedge X_2\wedge X_3,
\]

% manuscript lines / position: 4–7
j'ai envie de prouver que c'est \unsure{un|une} hom \uline{constant} i.e. qu'il se factorise par une \uline{section} de $X_1\wedge X_2\wedge X_3$, ce qui signifie aussi que l'hom composé précédent

% manuscript lines / position: 8–9
\diagram{\begin{tikzcd}[ampersand replacement=\&,column sep=small,row sep=small]
 B \arrow[rr,"f"] \arrow[dr] \& \& A\\
 \& \left(\bigotimes A_i\right)^{\unsure{\gamma|\mathcal G|\Gamma}}\arrow[ur] \&
\end{tikzcd}}

% manuscript lines / position: 10–13
est à valeurs dans $k.1_A$. (Du moins, \unsure{il y équivaut|il y a équivalence|on y aboutira} quand on suppose $k\unsure{\mathord{-}|\longrightarrow}A$ injectif --- c'est ce qu'on va supposer maintenant, pour simplifier.)

% manuscript lines / position: 14
Mais on a évidemment

% manuscript lines / position: 15
\[
 f(B)\subset A^\gamma,
\]

% manuscript lines / position: 16
et la condition est vérifiée si on suppose

% manuscript lines / position: 17
\[
 \boxed{A^\gamma=k.\unsure{1|\bar 1}_A}.
\]

% manuscript lines / position: 18–21
(C'est là une condition pas tjrs vérifiée, \unsure{$=$|même} dans des cas intéressants, à cause de la car. 2. Mais si 2 est inversible \deleted{\unsure{,}} dans $k$, alors tjrs vrai \unsure{compte tenu|au terme} des conditions \unsure{précdtes|précédentes}).

% manuscript lines / position: bottom left
% archive sheet number: 117
